\section{The joint large-distance tail of Group II Row IV: term-by-term check}\label{sec:tailfix}
\noindent\textbf{Result in one paragraph.} In the region where the measured gluon, the unmeasured gluon and their parent are all far from the projectile together ($w,w',z,y'\to\infty$ at fixed $w-z$, $w'-z$, $y'-z$; ``joint limit''), the finite terms \emph{as written} in \texttt{Final\_terms\_checked} / \texttt{Finite\_terms\_final} carry the tails, in units
\begin{align}
\mathsf{u}\equiv2N_{c}\,\mathcal{C}\,\frac{4\pi}{k^{2}}\int\frac{d^{2}Y}{Y^{2}}\,G_{\infty}(Y),\qquad G_{\infty}(Y)=\int_{x,x'}\big[U(Y)-U(x')\big]^{ce'}\big[U(Y)-U(x)\big]^{ce}\rho^{e'}(x')\rho^{e}(x),\qquad\varphi=\angle(k,Y),
\label{ft:unit}
\end{align}
\begin{center}
\renewcommand{\arraystretch}{1.3}
\begin{tabular}{|l|c|}\hline
Group II Row IV: DGLAP $+$ RC $+$ finite (remnant, tensor term) & $-\frac{11}{6}+\frac16\cos2\varphi$\\
Group II Row I, second part ($N_{2a}+\dots+N_{2f}$, i.e.\ $\sum\mathfrak{K}_{2\cdot}$) & $-\frac{11}{6}+\frac16\cos2\varphi$\\
Group III Row III, second part, with c.c.\ ($N_{2}^{\rm reg}+N_{2}^{\Psi}+N_{2}^{\rm sub}+N_{2}^{\log}$) & $+\frac{11}{3}-\frac13\cos2\varphi$\\
Group II Row I first part; Group III Row III first part & $0$\\\hline
sum & $0$\\\hline
\end{tabular}
\end{center}
and the $\log\frac{\vee}{k^{+}}$ (BFKL) and $\log\frac{k^{+}}{\Lambda}$ (JIMWLK) coefficients of the three rows carry $1+1-2=0$ times $\mathsf{u}$ each. \textbf{The tail cancels; no correction of the written finite terms is needed.} The scaleless rule and the recoil rewritings do move tail pieces between terms, but only inside the finite sector of Group III Row III (the term $N_{2}^{\log}$ and the extra term $N_{2}^{\Psi}$ each carry $+2\mathsf{u}$; the subtracted term $N_{2}^{\rm sub}$ carries none) and between the BFKL list and $N_{2}^{\log}$; the bookkeeping is complete in each sector. This holds for any target that is uniform at large distance ($U\to V$, e.g.\ a finite target, $V=1$), and for a colour-neutral projectile for \emph{any} target; for a coloured projectile and a target that is not uniform at large distance a remainder proportional to the total charge is left (the non-singlet term already listed as OPEN~II/III-4(b)). What changes are \emph{statements} of \texttt{Finite\_terms\_final} (red items in Sec.~\ref{ft:corr}). Scripts in \texttt{fix/tail/}: \texttt{t1\_density.py} (amplitude Fourier transforms, tail densities), \texttt{t2\_colour.py} (colour factors, $SU(2)$, $SU(3)$), \texttt{t3\_sectors.py} (sector integrals, Row~III split), \texttt{t4\_rowI\_written.py} (the written kernels of Row~I at large distance, 60 digits); outputs \texttt{t*.out}. [A]: analytic (sympy); [N]: numerical.

\subsection{Setup}
\noindent Phase $e^{-ik\cdot(w-w')}$, $k=(|k|,0)$, $t\equiv p^{+}/k^{+}$, $\xi=\frac{1}{1+t}$; $\mathcal{C}=\frac{1}{(2\pi)^{3}}\frac{g^{4}}{16\pi^{4}}\frac{1}{k^{+}}$ as in Sec.~\ref{sec:poles-dglap}. Joint limit: $x,x'$ fixed, $z=Zn$, $w=z+a$, $w'=z+a'$ (Group II) resp.\ $y'=z+b$ (Group III, conjugate amplitude after the $w'$ integration), $n=(\cos\varphi,\sin\varphi)$, $Z\to\infty$ at fixed $a,a',b$. In this limit every two-gluon amplitude is ${\cal O}(1/Z)$ with a coefficient of the form $c\cdot a/a^{2}$ (the $1/q$ Weizs\"acker--Williams pole of the soft parent, Sec.~\ref{sec:poles-dglap}, Step 14), so each row is $\frac{1}{Z^{2}}\times$(a function of $a,a'$ or $a,b$), whose Fourier transform is elementary: $\int d^{2}a\,e^{-ik\cdot a}\frac{a^{i}}{a^{2}}=-2\pi i\frac{k^{i}}{k^{2}}$, and for the conjugate $C$ amplitude of Group III $\int d^{2}b\,e^{i(1+t)k\cdot b}\frac{b^{m}}{b^{2}}=\frac{2\pi i}{1+t}\frac{k^{m}}{k^{2}}$. $\int d^{2}Z/Z^{2}$ is the logarithmic tail; ${\cal O}(1/Z^{3})$ and terms odd in $n$ at order $1/Z$ do not contribute. The expansion is uniform in $t$ (the scales in $t$ are set by $(x-w)^{2}/(x-z)^{2}\to1$), so the $p^{+}$ integral and the joint limit commute; this is checked on the written Row~I kernels in Step~4.

\noindent\textbf{Colour} [N, \texttt{t2}]. With the far Wilson lines equal, $U(w)=U(w')=U(z)=U(y)=U(y')=V$, the colour factors of the three rows, as written, are exactly $+N_{c}\Gamma$ (Group~II Row~I, all of $N_{2a}\dots N_{2f}$), $+N_{c}\Gamma$ (Group~II Row~IV before the reduction), $-N_{c}\Gamma$ (Group~III Row~III, second part), $\Gamma=[V-U(x')]^{ce'}[V-U(x)]^{ce}\rho^{e'}(x')\rho^{e}(x)$: deviations $\le9\times10^{-15}$ for $|\Gamma|\simeq0.5$--$2.7$, three random configurations each, $SU(2)$ and $SU(3)$.

\subsection{Step 1: the pointwise tail density}
\noindent [A, \texttt{t1}] With the amplitudes as written (Group~II Row~I: the bracket of the checked row before the $p^{+}$ integral, prefactor $\frac{g^{4}}{4\pi^{4}}\frac{dp^{+}}{2\pi}\frac{p^{+}k^{+}}{(p^{+}+k^{+})^{2}}\frac{1}{DD'}$; Row~IV: $\mf{C}_{1}\mf{A}^{(1)}$ with $y=z+\xi(w-z)$, prefactor $\frac{p^{+}k^{+}}{(p^{+}+k^{+})^{4}}$; Group~III Row~III, second part: the bracket of \texttt{Final\_terms\_checked} with $\frac{1}{p^{+}}=[\frac{k^{+}}{p^{+}(p^{+}+k^{+})}]_{+}+\frac{1}{p^{+}+k^{+}}$ before the subtraction and the $-\frac{p^{+}}{k^{+2}}$ term before its recoil subtraction), the Fourier coefficients of the $1/Z$ amplitudes of $\mf{B}_{2}$ and of $\mf{C}_{1}\mf{A}^{(1)}$ differ only by the normalisation factor $1+t$ of the two conventions, e.g.\ $(i,j)=(1,1)$: $-2\pi i\frac{(t^{2}+t+1)\cos\varphi}{t(t+1)}$ against $-2\pi i\frac{(t^{2}+t+1)\cos\varphi}{t}$. The three rows have the same tail density per $dt$ (in units $\mathcal{C}\,\Gamma/Z^{2}$, colour stripped, $|k|=k^{+}=1$), $8\pi\,\varrho(t)$,
\begin{align}
\varrho(t)=\frac{(t^{2}+t+1)^{2}+t^{2}\cos2\varphi}{t\,(1+t)^{4}}\,,\qquad \varrho(t)\,dt=\Big[\frac{1-\xi}{\xi}+\frac{\xi}{1-\xi}+\xi(1-\xi)\big(1+\cos2\varphi\big)\Big]d\xi\,,
\label{ft:rho}
\end{align}
so that with the colour factors, and the c.c.\ of Group~III (the density is real and even in $k$), the weights are $+1$, $+1$, $-2$ and the sum vanishes pointwise in $t$. In units $\mathsf{u}$ the density of each row is $(+1,+1,-2)\times\varrho(t)\,dt$. This is the coordinate-space form of the amplitude identity $b-cA=-\frac{V_{\rho g}V_{ggg}}{E_{2}(E_{2}-E_{1})}$.

\subsection{Step 2: Group II Row IV (reference)}
\noindent [A, \texttt{t3} [4]] $\int_{1/V}^{1/(1+\lambda)}\varrho\,d\xi=\log V+\log\frac1\lambda-\frac{11}{6}+\frac16\cos2\varphi$ ($V=\vee/k^{+}$, $\lambda=\Lambda/k^{+}$). Written: BFKL $\log V$ (pure log, Final\_terms form), JIMWLK $\log\frac1\lambda$ (coordinate form, tail from $\log(y'-y)^{2}$), DGLAP $+\frac1{12}$, RC $-\frac{11}{12}$, finite: remnant $-1$, tensor $+\frac16\cos2\varphi$. [In the form of \texttt{Finite\_terms\_final}, where $[\frac{1-\xi}{\xi}]_{+}$ acts on $(1-\xi)\mathcal{F}$, the remnant is inside DGLAP: $-\frac{11}{12}$, $-\frac{11}{12}$, $+\frac16\cos2\varphi$; same sum.]

\subsection{Step 3: Group II Row I, second part}
\noindent The written finite term is $\lim\big[\int_{\lambda}^{V-1}dt\,I(t)-c_{B}\log V-c_{J}\log\frac1\lambda\big]$, $c_{B}=\frac{1}{BB'}[\text{bracket of }\mathfrak{K}_{2c}]$, $c_{J}=\frac{1}{AA'}[\text{bracket of }\mathfrak{K}_{2b}]$ (Finite\_terms\_final, grey note), i.e.\ $\int_{0}^{1}dt\,(I-\frac{c_{J}}{t})+\int_{1}^{\infty}dt\,(I-\frac{c_{B}}{t})$. In the joint limit $c_{J},c_{B}\to\lim_{t\to0,\infty}t\varrho=1$ and [A, \texttt{t3} [2]]
\begin{align}
\text{tail}\Big(\sum\mathfrak{K}_{2\cdot}\Big)=\int_{0}^{\infty}dt\,\frac{t\cos2\varphi-(2t^{2}+3t+2)}{(1+t)^{4}}\;\mathsf{u}=\Big(-\frac{11}{6}+\frac16\cos2\varphi\Big)\mathsf{u}\,,\qquad \text{BFKL: }+\log V\,\mathsf{u},\quad\text{JIMWLK: }+\log\tfrac1\lambda\,\mathsf{u}.
\label{ft:rowI}
\end{align}
\textbf{Step 4: the written kernels} [N, \texttt{t4}, 60 digits]. $Z^{2}\sum_{2a}^{2f}\mathfrak{K}$, evaluated from the formulas as written (with $d_{\perp}=2$ and the $-1$ of $\mathfrak{K}_{2f}$; $x=(0.3,-0.2)$, $x'=(-0.1,0.4)$, $n$ and $-n$ averaged), against the leading product form $L(a,a')=\int dt\,[\frac{t}{(1+t)^{2}}T_{1}(a)\!\cdot\!T_{1}(a')-\text{subtractions}]$ whose Fourier transform gives \eqref{ft:rowI}: three random $(a,a',n)$, values $-1.99535$, $1.97289$, $-4.28034$; differences $1.3\times10^{-6},\,1.3\times10^{-8},\,1.3\times10^{-10}$ (resp.\ $4.8\times10^{-7\dots-11}$, $1.0\times10^{-6\dots-10}$) at $Z=10^{3},10^{4},10^{5}$. The single $\mathfrak{K}_{2f}$ is ${\cal O}(Z^{-2})$ as well (its ${\cal O}(Z^{0})$ parts, $\frac{d_{\perp}}{4CC'}\times$[$-\frac{n\cdot a'}{n\cdot(a'-a)}$, $-1$, $-\frac{n\cdot a}{n\cdot(a'-a)}$ with signs], cancel). The written form also equals the subtracted $p^{+}$ integral at these large-distance points to all digits shown. Since the $x,x'$ dependence enters only at ${\cal O}(1/Z^{3})$, the check with $x\ne x'$ confirms that the tail is $\propto G_{\infty}$.

\subsection{Step 5: Group III Row III, second part, as written}
\noindent The written second part (\texttt{Final\_terms\_checked}) has three terms: T1 (main bracket, plus prescription at $p^{+}\to0$), T2 (the $-\frac{p^{+}}{k^{+2}}$ term with the recoil subtraction $\mathcal{K}(x'-y')\mathcal{L}(y')-\mathcal{K}(x'-z)\mathcal{L}(z)$), T3 (the term with $\mathcal{K}(w'-z)\mathcal{K}(x'-z)\frac{(x-w)^{2}}{(x-z)^{2}(p^{+}+k^{+})D}$). T3 is what is left of the added-back subtraction of T2 after $-\log\frac{\vee}{k^{+}}\frac{1}{k^{+}(x-z)^{2}}\mathcal{W}\,\mathcal{K}\mathcal{K}\,\mathcal{L}(z)$ went to the BFKL list: with $\int_{y'}e^{i(1+t)k\cdot(y'-z)}\mathcal{K}(y'-z)=\frac{1}{1+t}\int_{w'}e^{ik\cdot(w'-z)}\mathcal{K}(w'-z)$ and $-\frac{p^{+}}{k^{+}(p^{+}+k^{+})D}=-\frac{1}{k^{+}B(p^{+}+k^{+})}+\frac{A}{B(p^{+}+k^{+})D}$. Per unit colour weight (multiply by $-2$) [A, \texttt{t3} [1], [3]]:
\begin{center}
\renewcommand{\arraystretch}{1.35}
\begin{tabular}{|l|l|c|c|}\hline
term & density in the joint limit & tail (finite) & $\times(-2)$\\\hline
T1 (main term), author's $[\,]_{+}$ & $\varrho_{A}=\frac{t^{2}\cos2\varphi+2t^{2}+2t+1}{t(1+t)^{4}}$, minus $\frac{1}{t(1+t)}$ & $-\frac56+\frac16\cos2\varphi$ & $\frac53-\frac13\cos2\varphi$\\
\quad $=N_{2}^{\rm reg}$ (phase kept in the subtraction) & minus $\frac{1}{t(1+t)^{2}}$ & $\frac16+\frac16\cos2\varphi$ & $-\frac13-\frac13\cos2\varphi$\\
\quad $+\,N_{2}^{\Psi}=f(0)\Psi$ & $\frac{1}{t(1+t)^{2}}-\frac{1}{t(1+t)}$ & $-1$ & $+2$\\
T2 $=N_{2}^{\rm sub}$ & $\mathcal{K}(x'-y')-\mathcal{K}(x'-z)={\cal O}(b/Z^{2})$ & $0$ & $0$\\
T3 $=N_{2}^{\log}$ & $-\frac{1}{(1+t)^{2}}$ & $-1$ & $+2$\\\hline
BFKL list (moved) & $\frac{1}{1+t}$ over $(0,V-1)$ & $\log V$ & $-2\log V$\\
JIMWLK list & $c_{J}=\lim t\varrho_{A}=1$ & $\log\frac1\lambda$ & $-2\log\frac1\lambda$\\\hline
\end{tabular}
\end{center}
Checks: $\varrho_{A}+\frac{t}{(1+t)^{2}}=\varrho$ (T1 plus the unsubtracted T2 is the whole row); T3 $+$ BFKL-moved $=\frac{t}{(1+t)^{2}}$, the density of the unsubtracted T2 (the recoil rewriting is exact). Total of the written finite part of Row~III with c.c.: $\big(\frac{11}{3}-\frac13\cos2\varphi\big)\mathsf{u}$. T1 has no $\log\vee$: $\varrho_{A}={\cal O}(t^{-3})$. The first parts of Group~II Row~I and Group~III Row~III have no joint tail: there the measured gluon of one side enters only through $\mathcal{K}(x'-w')$ (resp.\ $\mathcal{K}(y-w)$) and the phase, and $\int_{w}e^{-ik\cdot w}\mathcal{K}(y-w)$ ties it to the projectile [A]; their large-distance region is the soft one (\texttt{fix\_soft}).

\subsection{Step 6: sums by sector}
\noindent [A, \texttt{t3} [5]] Finite $+$ DGLAP $+$ RC: $\big(-\frac{11}{6}+\frac16\cos2\varphi\big)+\big(-\frac{11}{6}+\frac16\cos2\varphi\big)+\big(\frac{11}{3}-\frac13\cos2\varphi\big)=0$. BFKL: $(1+1-2)\log V=0$; JIMWLK: $(1+1-2)\log\frac1\lambda=0$. The sector-wise cancellation is not automatic (the three rows use different end-point subtractions: $\frac{1}{t}$ split at $t=1$ in Row~I, $\frac{1}{t(1+t)}$ and $\frac{1}{1+t}$ in Row~III, $\frac{1}{\xi},\frac{1}{1-\xi}$ in Row~IV); it holds because every row's split is an exact rewriting over the same range $\Lambda<p^{+}<\vee-k^{+}$ and the end-point coefficients are equal ($1,1,-2$) at both ends. The BFKL and JIMWLK \emph{formulas} are consistent with this: the only BFKL block with all of $w,w',z$ far is $\mathcal{K}(w-z)\cdot\mathcal{K}(w'-z)\,\mathcal{K}(x-z)\cdot\mathcal{K}(x'-z)$, whose colour factor vanishes for uniform far lines up to the commutator (contact) term [N, \texttt{t2}: $P_{1}+P_{2}\le1.2\times10^{-14}$] and contains no term bilinear in $U(x),U(x')$; the JIMWLK formula is $-H_{\rm JIMWLK}$ on the LO structure, and $(L_{w}-U(z)R_{w})U(w)=0$ for $U(z)=U(w)$, while $H$ does not act on the $U(x')U(x)$ part. So neither formula carries a joint tail; the $\pm\log$ tails of the single rows cancel inside them.

\noindent\textbf{Where the scaleless rule enters.} Nowhere in this region: the dropped poles $\frac{2}{\epsilon}$ of Row~IV multiply $\int d^{2}u\,e^{ik\cdot u}G_{\infty}/Y^{2}\propto\delta^{(2)}(k)$ (no tail at $k\ne0$); in Row~I no pole multiplies $d_{\perp}$ (the ${\cal O}(Z^{0})$ parts of $\mathfrak{K}_{2f}$ cancel at $d_{\perp}=2$, Step~4).

\subsection{Step 7: colour-neutral projectile and general target}
\noindent [N, \texttt{t2}] (a) Uniform far lines (finite target, $V=1$, or any constant $V$), coloured or not: the three colour factors are $N_{c}\Gamma,N_{c}\Gamma,-N_{c}\Gamma$ (Step~1): cancellation as above. (b) Colour-neutral projectile, arbitrary target ($U(w),U(w'),U(z),U(y),U(y')$ all different and depending on the relative coordinates): at order $1/Z^{2}$ the kernels do not depend on $x,x'$, so every part of a colour factor that is not bilinear in $U(x),U(x')$ multiplies $\int_{x}\rho(x)$ or $\int_{x'}\rho(x')$ and vanishes on a singlet state ($Q^{a}$ stands next to the state in the written order). The bilinear parts are far-line independent: $+N_{c}U^{T}(x')U(x)$, $+N_{c}U^{T}(x')U(x)$, $-N_{c}U^{T}(x')U(x)$ for Rows I, IV, III [deviations $\le3.3\times10^{-15}$, $SU(2)$, $SU(3)$]; T2 has no bilinear part in $\mathcal{L}(y')-\mathcal{L}(z)$. The tail cancels for any target. (c) Coloured projectile, target not uniform at large distance: the charge parts do not cancel (norm of the remainder of $I+IV+2\,III$: $3.6$ ($SU(2)$), $4.8$ ($SU(3)$), zero to $5\times10^{-15}$ with uniform lines); this is the term $[U(w)U(z)-U(y)]\,b\,Q$ of the amplitude (Sec.~\ref{sec:poles-dglap}, Step 14), the non-singlet large-distance term of OPEN~II/III-4(b), not a defect of the finite terms.

\subsection{Corrections to statements}\label{ft:corr}
\noindent No term of the written finite terms has to be changed. The following statements are corrected:
\begin{enumerate}
\item {\color{red}\texttt{Finite\_terms\_final}, OPEN~II/III-1, ``No term of Groups II, III cancels it'': the tail $2N_{c}\mathcal{C}\frac{4\pi}{k^{2}}\int\frac{d^{2}Y}{Y^{2}}[-\frac{11}{6}+\frac16\cos2\varphi]G_{\infty}$ of Group~II Row~IV (i)+(ii)+(iii) is cancelled by Group~II Row~I, second part ($\sum\mathfrak{K}_{2\cdot}$: $-\frac{11}{6}+\frac16\cos2\varphi$) and Group~III Row~III, second part with c.c.\ ($+\frac{11}{3}-\frac13\cos2\varphi$), in the region where $w,w',z$ (and $y'$) are far together. The terms are integrable at large distance only together: the $d^{2}Y$ (pair position) integral of these three rows has to be done under one integral, with $Y=\frac{y+y'}{2}$ of Row~IV identified with $z$ (or any point of the pair) of the other rows.}
\item {\color{red}\texttt{Finite\_terms\_final}, Sec.~``Large-distance tails'', list of terms ``without a $1/z^{2}$ tail'': true at fixed $w,w'$ only. In the joint region $\sum\mathfrak{K}_{2\cdot}$, $N_{2}^{\rm reg}$, $N_{2}^{\log}$ do have $1/Z^{2}$ tails (values in Step~5); $N_{2}^{\Psi}$ has, besides its soft tail of the table, a joint tail $+2\mathsf{u}$ (with c.c.).}
\item {\color{red}Sec.~\ref{sec:poles-dglap}, ``What remains (item 4)'': done; the compensating tail is present in the written finite terms, sector by sector, and requires no change of the RC/DGLAP split. With the local (tail-free) DGLAP and RC of Step~12, the remainder $N_{\mathcal{R}}$ carries the whole Row~IV tail and must be grouped with $\sum\mathfrak{K}_{2\cdot}$ and Group~III Row~III (second part) as stated there.}
\end{enumerate}

\subsection{Not done / open}
\noindent (1) The joint limit was taken with $x,x'$ fixed and the far lines uniform (or the neutral-projectile reduction); the transition region where only some of $w,w',z$ are far (e.g.\ $|a|\sim Z$) was not analysed beyond the convergence of the Fourier integrals of the leading terms (oscillation at $|a|\gg1/|k|$). (2) The commutativity of the $p^{+}$ integral and the joint limit is verified numerically for Row~I only (Step~4); for Row~III the terms are $p^{+}$ integrands and the limit was taken at fixed $t$ with $\varrho_{A}={\cal O}(t^{-3})$, $\frac{1}{(1+t)^{2}}$ integrable. (3) No full transverse integration with Wilson lines. (4) The statement about the BFKL and JIMWLK formulas is analytic (colour of the one far block; $-H_{\rm JIMWLK}$ form as checked in \texttt{Final\_terms\_checked}), not a term-by-term numerical evaluation of the written formulas in the joint limit. (5) The coloured-projectile remainder of Step~7(c) is not addressed here.
